% GENERATED FILE. Edit canonical lessons, then run npm run generate:books.
% canonical-source-hash: fnv1a64:70439abbbed219f4
% canonical-lessons: LA-C501-commoveo, LA-C501-afficio, LA-C501-affectus, LA-C501-tranquillus, LA-C501-securus

\chapter{Moods and feelings}
\label{ch:la-moods-and-feelings}
\hlchaptermodality{501}

\begin{chapteropening}
\textbf{By the end of this chapter:} \emph{I can describe moods and how things affect me.}

Complete the last lesson of chapter 501: I can describe moods and how things affect me.
\end{chapteropening}

\section[commoveō, commovēre]{commoveō, commovēre — to move deeply, disturb}

\label{lesson:LA-C501-commoveo}



\begin{quote}
Before the new one: say the Latin for to envy, then the Latin for a rival.
\end{quote}

\subsection*{You'll want to know: commoveō, commovēre}
\textbf{commoveō, commovēre} — \textquotedblleft{}to move deeply, disturb\textquotedblright{}.

\begin{quote}
Lacrimae fīliae patrem valdē commovent.
\end{quote}

\begin{quote}
\emph{His daughter's tears move the father deeply.}
\end{quote}

\emph{commoveō, commovēre, commōvī, commōtum}. Of the body, to shake; of the mind, to stir or upset. English \emph{commotion}.

\begin{cousinweb}
\textbf{commoveō} is built on \emph{con-} and \emph{moveō} \textquotedblleft{}to move\textquotedblright{} (Lewis \& Short; EtymDB). English \emph{commotion} comes from the related noun \emph{commōtiō}.
\end{cousinweb}

\begin{grammarlens}[title={What you've built}]
One more word for telling, explaining and describing at B1.
\end{grammarlens}

\subsection*{Your turn}
\begin{itemize}
\raggedright
  \item \emph{Say it:} \emph{commoveō, commovēre}
  \item \emph{Say it:} \emph{commoveō, commovēre}, once more
  \item \emph{Say it:} a sentence of your own with \emph{commoveō, commovēre}
  \item \emph{Recall:} say the Latin for to envy, then the Latin for a rival, then say \emph{commoveō, commovēre} again
\end{itemize}

\begin{tcolorbox}[breakable,colback=teal!4,colframe=teal!35!black,title={Before you move on}]
What does commoveō, commovēre mean? (\textquotedblleft{}To move deeply, disturb\textquotedblright{}.) Say it once more.
\end{tcolorbox}
\section[afficiō, afficere]{afficiō, afficere — to affect, fill with a feeling}

\label{lesson:LA-C501-afficio}



\begin{quote}
Before the new one: say the Latin for a rival, then the Latin for to move deeply.
\end{quote}

\subsection*{You'll want to know: afficiō, afficere}
\textbf{afficiō, afficere} — \textquotedblleft{}to affect, fill with a feeling\textquotedblright{}.

\begin{quote}
Nūntius malus omnēs dolōre afficit.
\end{quote}

\begin{quote}
\emph{The bad news fills everyone with grief.}
\end{quote}

\emph{afficiō, afficere, affēcī, affectum}. \emph{Aliquem laetitiā afficere} — to fill someone with joy; in the passive, \emph{dolōre afficior} — I am stricken with pain. English \emph{affect}.

\begin{cousinweb}
\textbf{afficiō} is built on \emph{ad-} and \emph{faciō} \textquotedblleft{}to do, make\textquotedblright{} (Lewis \& Short): to do something to someone. English \emph{affect} comes from it.
\end{cousinweb}

\begin{grammarlens}[title={What you've built}]
One more word for telling, explaining and describing at B1.
\end{grammarlens}

\subsection*{Your turn}
\begin{itemize}
\raggedright
  \item \emph{Say it:} \emph{afficiō, afficere}
  \item \emph{Say it:} \emph{afficiō, afficere}, once more
  \item \emph{Say it:} a sentence of your own with \emph{afficiō, afficere}
  \item \emph{Recall:} say the Latin for a rival, then the Latin for to move deeply, then say \emph{afficiō, afficere} again
\end{itemize}

\begin{tcolorbox}[breakable,colback=teal!4,colframe=teal!35!black,title={Before you move on}]
What does afficiō, afficere mean? (\textquotedblleft{}To affect, fill with a feeling\textquotedblright{}.) Say it once more.
\end{tcolorbox}
\section[affectus, affectūs]{affectus, affectūs — a mood, state of mind}

\label{lesson:LA-C501-affectus}



\begin{quote}
Before the new one: say the Latin for to move deeply, then the Latin for to affect.
\end{quote}

\subsection*{You'll want to know: affectus, affectūs}
\textbf{affectus, affectūs} — \textquotedblleft{}a mood, state of mind\textquotedblright{}.

\begin{quote}
Affectūs animī in vōce audīmus.
\end{quote}

\begin{quote}
\emph{We hear the moods of the mind in the voice.}
\end{quote}

Masculine, fourth declension, from \emph{afficiō}. English \emph{affection} comes from the related \emph{affectiō}.

\begin{cousinweb}
\textbf{affectus} is the action noun of \emph{afficiō} \textquotedblleft{}to affect\textquotedblright{} (EtymDB; Lewis \& Short): the state into which something puts the mind.
\end{cousinweb}

\begin{grammarlens}[title={What you've built}]
One more word for telling, explaining and describing at B1.
\end{grammarlens}

\subsection*{Your turn}
\begin{itemize}
\raggedright
  \item \emph{Say it:} \emph{affectus, affectūs}
  \item \emph{Say it:} \emph{affectus, affectūs}, once more
  \item \emph{Say it:} a sentence of your own with \emph{affectus, affectūs}
  \item \emph{Recall:} say the Latin for to move deeply, then the Latin for to affect, then say \emph{affectus, affectūs} again
\end{itemize}

\begin{tcolorbox}[breakable,colback=teal!4,colframe=teal!35!black,title={Before you move on}]
What does affectus, affectūs mean? (\textquotedblleft{}A mood, state of mind\textquotedblright{}.) Say it once more.
\end{tcolorbox}
\section[tranquillus, tranquilla, tranquillum]{tranquillus, tranquilla, tranquillum — calm, peaceful}

\label{lesson:LA-C501-tranquillus}



\begin{quote}
Before the new one: say the Latin for to affect, then the Latin for a mood.
\end{quote}

\subsection*{You'll want to know: tranquillus, tranquilla, tranquillum}
\textbf{tranquillus, tranquilla, tranquillum} — \textquotedblleft{}calm, peaceful\textquotedblright{}.

\begin{quote}
Post tempestātem mare tranquillum est.
\end{quote}

\begin{quote}
\emph{After the storm the sea is calm.}
\end{quote}

Of the weather and the sea, and of a quiet mind. English \emph{tranquil}.

\begin{cousinweb}
\textbf{tranquillus} goes back to Proto-Italic (EtymDB), but its further origin is not settled. English \emph{tranquil} comes from it.
\end{cousinweb}

\begin{grammarlens}[title={What you've built}]
One more word for telling, explaining and describing at B1.
\end{grammarlens}

\subsection*{Your turn}
\begin{itemize}
\raggedright
  \item \emph{Say it:} \emph{tranquillus, tranquilla, tranquillum}
  \item \emph{Say it:} \emph{tranquillus, tranquilla, tranquillum}, once more
  \item \emph{Say it:} a sentence of your own with \emph{tranquillus, tranquilla, tranquillum}
  \item \emph{Recall:} say the Latin for to affect, then the Latin for a mood, then say \emph{tranquillus, tranquilla, tranquillum} again
\end{itemize}

\begin{tcolorbox}[breakable,colback=teal!4,colframe=teal!35!black,title={Before you move on}]
What does tranquillus, tranquilla, tranquillum mean? (\textquotedblleft{}Calm, peaceful\textquotedblright{}.) Say it once more.
\end{tcolorbox}
\section[sēcūrus, sēcūra, sēcūrum]{sēcūrus, sēcūra, sēcūrum — free from care, untroubled}

\label{lesson:LA-C501-securus}



\begin{quote}
Before the new one: say the Latin for a mood, then the Latin for calm.
\end{quote}

\subsection*{You'll want to know: sēcūrus, sēcūra, sēcūrum}
\textbf{sēcūrus, sēcūra, sēcūrum} — \textquotedblleft{}free from care, untroubled\textquotedblright{}.

\begin{quote}
Post labōrem agricola sēcūrus dormit.
\end{quote}

\begin{quote}
\emph{After his work the farmer sleeps free from care.}
\end{quote}

\emph{sē-} 'without' + \emph{cūra} 'care'. In classical Latin it means carefree rather than safe. English \emph{secure}.

\begin{cousinweb}
\textbf{sēcūrus} is built on \emph{sē-} \textquotedblleft{}without\textquotedblright{} and \emph{cūra} \textquotedblleft{}care\textquotedblright{} (Lewis \& Short; EtymDB). English \emph{secure} and \emph{sure} come from it.
\end{cousinweb}

\begin{grammarlens}[title={What you've built}]
One more word for telling, explaining and describing at B1.
\end{grammarlens}

\subsection*{Your turn}
\begin{itemize}
\raggedright
  \item \emph{Say it:} \emph{sēcūrus, sēcūra, sēcūrum}
  \item \emph{Say it:} \emph{sēcūrus, sēcūra, sēcūrum}, once more
  \item \emph{Say it:} a sentence of your own with \emph{sēcūrus, sēcūra, sēcūrum}
  \item \emph{Recall:} say the Latin for a mood, then the Latin for calm, then say \emph{sēcūrus, sēcūra, sēcūrum} again
\end{itemize}

\begin{tcolorbox}[breakable,colback=teal!4,colframe=teal!35!black,title={Before you move on}]
What does sēcūrus, sēcūra, sēcūrum mean? (\textquotedblleft{}Free from care, untroubled\textquotedblright{}.) Say it once more.
\end{tcolorbox}
